\section{Conclusion}
\label{sec:conclusion}

The 2024 to 2026 wave of inter-agent protocols is the
third generation of an agent-communication-language
lineage that began with KQML and FIPA-ACL in the
mid-1990s. The contemporary wave differs from its
predecessors in two structural ways: the receivers
of inter-agent messages are LLMs with a flexible
natural-language interlingua, which absorbs schema
mismatches the FIPA era could not handle; and the
protocols carry a third architectural layer
(commerce settlement) that no prior wave needed to
address. The flexibility is the source of the
contemporary protocols' rapid adoption and of their
distinctive threat surface; the third layer is the
source of the institutional traction the
agentic-commerce stack has accumulated in twelve
months.

This survey contributed seven distinct artifacts to the
literature on this wave
(\S\ref{sec:intro:contributions}). Contributions C1 to C4
are survey contributions. \textbf{C1}: a
uniform ten-axis rubric, applied to the eight
tool-binding and peer-messaging protocols of the thirteen
surveyed, rendered as
Table~\ref{tab:rubric-matrix},
Figure~\ref{fig:heatmap}, and Figure~\ref{fig:radar}.
\textbf{C2}: first-class treatment of the
agentic-commerce stack (AP2, Visa TAP, Mastercard
Agent Pay, ERC-8004, UCP) as a third architectural layer
of inter-agent coordination, alongside tool-binding
(MCP and family) and peer-messaging (A2A and family).
\textbf{C3}: an industry-adoption map across four
verticals (enterprise SaaS, agentic commerce and
payments, healthcare, finance) with named production
deployments, rendered as
Table~\ref{tab:industry-adoption}. \textbf{C4}: a
machine-readable companion artifact
(\texttt{data/protocol-scores.yaml},
\texttt{data/threat-coverage.yaml},
\texttt{data/industry-adoption.yaml}, and figure
regeneration scripts) that lets the rubric scores
remain re-runnable and updatable as protocol
versions drift across the 6 to 9 month review cycle.
The artifact makes each score traceable and each figure
regenerable. It does not stand in for a second rater,
which the survey does not have
(\S\ref{sec:rubric:calibration}).

Contributions C5, C6, and C7 are protocol-design
proposals derived from the open-problems agenda the
survey identifies, presented in \S\ref{sec:proposals}.
\textbf{C5} is the OWA scoping scheme: a three-tier
\emph{org / workspace / agent} URI extending the
\texttt{agent://} scheme of
Rodriguez~\cite{agent_uri_scheme} with an explicit workspace
tier that no contemporary identifier scheme provides.
\textbf{C6} is the AGRP guardrails envelope: a normative
extension to A2A wrapping every cross-boundary message
with classification labels, a fingerprint-evidenced
redaction log on the spans the egress filter touched,
and cryptographic attestation that the filter ran.
The envelope is partitioned into two normative assurance
tiers. A receiver can then distinguish a sender's claim
that a filter ran (Tier~1) from evidence signed by a
boundary the sender does not control (Tier~2). An
optional log-anchored sub-tier covers adversarial-audit
contexts.
\textbf{C7} is the ACSP context selection envelope: a
sibling extension to A2A that records the sender's
declared context-selection step itself: which candidate
items were considered, which were kept, which were
dropped, under which token budget and approximation
objective, with a signed wire-level audit log whose
observable invariants the receiver can verify, anchored on the PACMS reference
implementation~\cite{pacms_cikm_demo}. The three
proposals are co-designed: OWA names the boundary, ACSP
selects which items cross it, AGRP redacts the spans
within those items. The triplet closes a single loop of
naming, selection, and attestation. No pair
closes that loop individually.

The six principal findings of the survey are
worth restating in compact form.

\paragraph{Finding 1: The contemporary protocols
specialise rather than compete.}
The most-debated apparent rivalry in the literature
is MCP versus A2A. It is a category error
(\S\ref{sec:taxonomy:complementary}). The two
protocols occupy different architectural layers of
the three-layer taxonomy and compose in production
(\S\ref{sec:adoption:saas}) rather than
substituting. The same compositional logic extends
through the commerce-settlement layer and through
the ACP, AGNTCY, ANP, and Coral peer-messaging
alternatives.

\paragraph{Finding 2: No protocol scores uniformly
high on the rubric.}
The rubric matrix
(Table~\ref{tab:rubric-matrix}) renders an
observation the prior surveys could not make
precise: the contemporary protocols specialise
into mutually-exclusive strengths, and no surveyed
protocol combines a deployed task model with a
decentralised identity model. The single
composition the rubric most clearly recommends is
ANP's identity layer beneath A2A's task layer
(\S\ref{sec:comparison:stack}). That stack is not yet
deployed. It would deliver more capability
per spec-word than any other near-term composition.
The recommendation is for the identity layer specifically,
not for ANP as a whole: its DID mechanism carries no
dependency on the rest of the stack, while its description
and discovery protocols overlap the A2A AgentCard and
compete with it rather than stacking beneath it.

\paragraph{Finding 3: The authorisation-propagation
gap is the central open problem.}
Not one of the eight scored protocols specifies a
cryptographic delegation chain. The gap is universal, and
it survives the strongest counterexample the survey could
find. ACNBP carries the highest threat-coverage count in
the set and the most rigorous single-hop authorization
primitive, and it still binds only two parties at a time,
which a lead author confirmed is the intended scope. Every
surveyed protocol therefore leaves the delegation
mechanism to the surrounding identity architecture.
Deployments fill the gap with token
forwarding, which the security
literature~\cite{authz_propagation, agentic_jwt}
identifies as the dominant attack surface in
production. The gap is therefore an omission at the
specification layer with a measured cost at the
deployment layer. Rigour on one hop is not the same
property as authority that travels, and the survey's
Axis 3 column separates the two. The gap is the highest-priority item
on the open-problems agenda
(\S\ref{sec:openproblems:priorities}, Tier 1) and
the single change to the contemporary protocols
that would most materially improve the threat-coverage
profile of
Table~\ref{tab:threat-surface}.

\paragraph{Finding 4: The semantic-interoperability
equilibrium is local and unsustainable.}
The level-1 cluster on Axis 9
(\S\ref{sec:semantic}) is the contemporary
equilibrium and is convenient for protocol authors,
but the schema-fragmentation cost
(\S\ref{sec:semantic:fragmentation}) is visible in
the substantially lower inter-organisational call
success rate documented in AAIF working group
discussions~\cite{lf_aaif} and in the AgentLeak
benchmark's 68.8\% inter-agent leakage finding for
content forwarded across organisational
boundaries~\cite{agentleak}. The field has not
decided how to escape the equilibrium: ontology-base
or semantic-matchmaker, protocol-layer or
orchestrator-layer. The
answer is the single largest
research-and-standardisation question on the
agenda (\S\ref{sec:openproblems:schema}).

\paragraph{Finding 5: Governance, not technical
rigour, predicts adoption.}
\label{sec:conclusion:governance}
The protocols with the highest rubric scores
(ACNBP, LOKA) have the smallest deployment
footprints; the protocols with the largest
footprints (A2A, MCP) have only modest rubric
profiles. The pattern is consistent across the
adoption data (\S\ref{sec:adoption:synthesis}) and
the governance analysis (\S\ref{sec:governance}):
institutional positioning under foundations and
standards bodies drives adoption more strongly
than technical quality at the specification level.
The lesson for the research-and-standardisation
community is that the highest-impact technical
contributions will be those that institutional
incentives are positioned to adopt. That is an argument
for working within the AAIF and W3C tracks rather
than alongside them.

\paragraph{Finding 6: Naming, selection, and guardrails
are co-design problems, not orthogonal ones.}
The three proposals introduced in \S\ref{sec:proposals}
illustrate a structural observation that the survey
treats as the sixth finding: the workspace-naming gap
(Contribution C5), the context-selection gap
(Contribution C7), and the egress-attestation gap
(Contribution C6) are not independent open problems but
three legs of the same loop. A guardrails envelope
without a normative workspace name needs either a
duplicate naming standard or an inferred boundary
notion that loses cross-vendor verifiability; a
selection envelope without a normative boundary kind
has nothing to key its policy on and would re-derive
OWA implicitly; a naming scheme without selection and
guardrails provides scope but no protocol-layer
enforcement of the boundaries it identifies; and a
selection-and-guardrails pair without a naming layer
leaves the boundary detection to convention. The
OWA~+~ACSP~+~AGRP triplet closes the loop, and the
closure is the survey's clearest illustration of why
the rubric axes (\S\ref{sec:rubric}) cannot be
addressed independently: identity (Axes 1 to 2),
authorisation (Axis 3), and security (Axis 8) are
deeply coupled in any deployment that crosses an
organisational boundary, and the protocol revisions
that the field adopts will succeed or fail on whether
they treat the coupling as a first-class design
constraint.

\paragraph{What the next cycle should produce.}
The open-problems agenda
(\S\ref{sec:openproblems:priorities}) ranks the
P-items by tractability and impact. The Tier-1
items are the cryptographic delegation chain and
the ANP-to-A2A identity bridge. Both are
near-term tractable and high-impact. Both have
clear institutional homes in the AAIF and W3C
working groups. The Tier-2 items (the schema
fragmentation problem, cross-network mandate
harmonisation, cross-tenant policy enforcement)
will require multi-cycle work; the Tier-3 items
(prompt-injection mitigation as a spec requirement,
contextual integrity for forwarding, identity
rotation) require research before they can be
standardised; the Tier-4 methodological items
(rubric-stable benchmarking, multi-tenant
evaluation harness) are the instrument the field
needs to measure progress on the rest. Together
the items constitute what the survey treats as
the agenda for the 2026 to 2027 revision cycle.

\paragraph{The artifact this survey leaves behind.}
\label{sec:conclusion:artifact}
The rubric of \S\ref{sec:rubric} is intended to
outlast the specific protocol versions it scores.
In the companion repository
(\companionrepo),
\texttt{protocol-scores.yaml},
\texttt{threat-coverage.yaml}, and
\texttt{industry-adoption.yaml} carry the
machine-readable form of the rubric scores; the
figure-regeneration scripts produce the heatmap and
radar charts mechanically from those scores.
Future readers facing protocols that did not exist
at the 2026-Q1 cutoff can fork the YAML, extend the
rubric to the new protocols, regenerate the figures,
and produce a comparison that is directly comparable
to the comparison rendered here. The survey's
durable contribution is the rubric and the
discipline of using it; the specific scores will
drift as the protocols evolve, but the comparison
surface will remain. The OWA URI scheme, the AGRP
guardrails envelope, and the ACSP context selection
envelope (Contributions C5, C6, and C7) are part of
the same companion artefact tradition: each is
specified as a schema sketch with machine-readable
ABNF and JSON examples in \S\ref{sec:proposals}, so
that an implementer willing to adopt any of the
three proposals has a normative reference rather
than a paraphrase to work from.

The protocols of the 2024 to 2026 wave are still
forming, and the open-problems agenda is large
enough that the wave's eventual shape is not yet
legible. What is legible is that the field is
moving. It moves quickly enough that any survey of
inter-agent protocols has a sell-by date measured
in quarters. Also legible is that the rubric this survey
contributes is the instrument by which the field
can measure its own movement. We close in the
hope that the next survey of this landscape,
written eighteen months hence, finds the
P-items of \S\ref{sec:openproblems} substantially
closed, the rubric matrix's strongest cells
populated across more protocols, and the
agentic-AI ecosystem better-served by the
inter-agent protocols its participants are
collectively building.

\begin{acks}
The authors thank Joshua Waldrep, maintainer
of the Pipelock project~\cite{pipelock}, for a detailed
external review of the Pipelock description in
\S\ref{sec:proposals:agrp}. He pressed the
self-attestation vs.\ independent-attestation
distinction that shaped the two-tier assurance model of
\S\ref{sec:proposals:agrp:tiers}. That review also
produced two tightenings: independence must cover the
enforcement, not only the signature, and the verifier key
must be published out-of-band ahead of the receipt. The
optional log-anchored sub-tier came out of the same
exchange. Any remaining errors in
the description of Pipelock or in the tier definitions
are the authors' own.

The authors thank the ANP Community for the explanation of
ANP's design intent that shaped
\S\ref{sec:industry:anp}, credited at the community's own
request rather than to an individual. That exchange
established that Information and Interface, not tasks, are
ANP's first-class concepts, that a task is treated as one
form of an Interface on a generality argument, and that
asynchrony is defined per-Interface rather than as a
protocol primitive. It also corrected the survey's
composition recommendation: the portable layer is ANP's
identity mechanism, while Agent Description and Agent
Discovery overlap the A2A AgentCard and compete with it
rather than stacking beneath it. The Axis 5 and Axis 6
scores were not changed by the exchange, and the reading of
what those scores mean was.

The authors thank Ken Huang, lead author of
ACNBP~\cite{spec_acnbp}, who reviewed this survey's scoring
of that protocol across four exchanges. He is credited at his
own request and in the form he chose. The survey had scored
ACNBP at level 3 on Axis 3, making it the only protocol in
the set to reach that level and the sole counterexample to
Finding 3. He corrected the reading: ACNBP is a two-party
negotiation and binding protocol, and multi-step delegation
is not within its scope. The score moved to 2, and Finding 3
became the stronger claim that the authorisation-propagation
gap is universal across all eight scored protocols. The same
exchange established that the companion zero-trust identity
framework~\cite{huang_zerotrust_identity} is related work
rather than a normative dependency, which settles Axis 2,
and confirmed the threat-coverage count of
Table~\ref{tab:threat-surface}. His review covered the ACNBP
scores. It did not cover the rubric, the other seven scored
protocols, or the proposals of
\S\ref{sec:proposals}.

The authors thank Akram Sheriff for two rounds of review of
the full manuscript. Four parts of this survey exist because
he pressed for them. First, he argued that a flat count over
the twelve threat classes equates a low-severity mitigation
with a high-severity one. Residual risk, he noted, is the
quantity a deployer actually asks about. The severity-banded
weights and Table~\ref{tab:severity-coverage} answer that
objection, and \S\ref{sec:rubric:axis8} states the banding
criteria he asked us to make explicit. Second, he pressed
twice on scoring subjectivity. That produced the calibration
rules, the worked examples, and the plain statement of what
agreement evidence we do not have
(\S\ref{sec:rubric:calibration}). Third, he judged the three
proposals underspecified on their failure modes. That
produced the misclassification analysis of AGRP
(\S\ref{sec:proposals:agrp:open}, Q4) and the candidate-set
disclosure analysis of ACSP
(\S\ref{sec:proposals:acsp:open}, Q3). Fourth, he pointed us
at two works we had missed. The proof-of-guardrail
scheme~\cite{proof_of_guardrail} prompted the separation of
what AGRP guarantees with and without a trusted execution
environment (\S\ref{sec:proposals:agrp:tee}).
AgentDNS~\cite{agentdns} prompted the trust-root contrast of
\S\ref{sec:proposals:owa:related}. His review did
not extend to the individual protocol scores.

We also contacted Shopify, Microsoft, and Cisco about
specific claims on protocols they author or govern. Each
response changed the text, and the body says where. Shopify
confirmed the orchestration-layer reading of UCP and
corrected its treatment of request signing and identity
linking (\S\ref{sec:emerging:commerce}). Microsoft corrected the
token-forwarding claim about A2A, named the extension route
for a delegation chain, and raised the denotation question of
\S\ref{sec:proposals:owa:open}. Cisco pointed us to A2A's
explicit exclusion of authorization semantics and argued that
ACNBP's breadth limits its adoption
(\S\ref{sec:governance}). None of the three reviewed the
survey as a whole.

No acknowledgement here implies
endorsement of the survey's assessment of any protocol.
Every score remains the authors' own judgement.
\end{acks}
